\section{The inner path and its algebraic representation}\label{sec:inner-algebra}

We construct the algebraic part of the decision gadget promised in
Theorem~\ref{thm:port-gadget}.  Its graph consists of one directed path for
each tuple of read-only inputs.  Besides encoding these paths, the
construction arranges that every chart involving one read port has exactly
the same native layout and formula in a common one-port template.  This
literal correspondence, including all verification auxiliaries, will supply
the distinguished source in Section~\ref{sec:inner-gadget}.

\emph{All field, clock, dimension, and size notation is reset in this
section.}  In particular, the fields and the polynomial-size affine domain
below are independent of the outer index group.  The construction and its
proof do not use the outer circuit realization.  All graphs in this section
are described implicitly; we never enumerate their exponentially many
Boolean input tuples.

\subsection{Parameters and the extended path}\label{ia:path-parameters}

Fix the exponent allowance $c_{\rm allow}$ in
Theorem~\ref{thm:port-gadget}, and choose a fixed integer
$c_0>c_{\rm allow}+5$.  For a port length $L$, set
\begin{equation}
 n'=(L+2)^{c_0},\qquad
 h=\Theta\bigl((\log(n'+2))^{1/4}\bigr),\qquad
 H=\mathbb F_h,\qquad q=h^v,                         \label{ia:parameters}
\end{equation}
where $h$ is a deterministically chosen power of two and $v$ is a fixed
sufficiently large integer.  Let $\ell$ be the least positive integer
coprime to $v$ for which $h^\ell\ge(n')^3+3$, and put
\begin{equation}
 K_1=\mathbb F_{h^\ell},\quad K=\mathbb F_{q^\ell},\quad
 C=K_1^\times,\quad M=h^\ell-1,\quad D=\ell^2-1.
                                                               \label{ia:fields}
\end{equation}
Fix a generator $\gamma$ of $C$.  Since $\gcd(\ell,v)=1$, the subfields
$K_1$ and $\mathbb F_q$ of $K$ intersect in $H$ and generate $K$.
An $H$-basis of $K_1$ is therefore an $\mathbb F_q$-basis of $K$.
We use it throughout: one $K$ variable has $\ell$ coordinates over
$\mathbb F_q$, and its $H^\ell$ digit grid is precisely $K_1$.

An interval of $v$ consecutive integers contains an integer congruent to
$1$ modulo $v$.  Thus the coprimality restriction changes the first size
threshold for $\ell$ by at most $v$, and
\begin{equation}
 \ell=\Theta(\log n'/\log\log n'),\qquad
 h=o(\ell),\qquad |K|=q^\ell=(n')^{O(1)}.          \label{ia:size-bounds}
\end{equation}
Choose $v$ large enough that $q$ dominates the fixed power of $D$ needed
by the scalar line--point input in Section~\ref{sec:inner-binary}, with
positive power slack.  In particular, $D/q\longrightarrow0$.  After $v$,
fix an integer $a>v$ and set
\begin{equation}
 B=\lceil q^{1/a}\rceil\le\ell                     \label{ia:radix-size}
\end{equation}
for all sufficiently large $L$.  Every fixed multiple of $\ell h$ is then
eventually below $D$.  All these choices precede the number of ports and
the particular Boolean function.

All finite sets used below have polynomial cardinality in $n'$.  Field
presentations, compatible embeddings, bases, generators, interpolation
matrices, and deterministic orderings can consequently be found by
exhaustive finite-field algebra in polynomial time in these cardinalities.
For instance, irreducible presentations and suitable generators can be
found by enumeration and their defining properties checked by finite
arithmetic.  No claim of polynomial time in the logarithm of a field's
cardinality is required.  Fix these conventions once for the common
template.  Finitely many smaller sizes can be handled by enlarging the
fixed lower cutoff as in the gadget statement.

Now fix an arbitrary constant number $k\ge1$ of ports and an arbitrary
constant size factor for the Boolean circuit of
$f:(\{0,1\}^L)^k\to\{0,1\}$.  Invocations with fewer ports are padded
by unused public-zero ports.  Normalize the circuit to a topological list
of gates with Boolean rules of degree at most two over characteristic two,
and finish with a separate output copy gate.  This costs only a constant
factor in circuit size.  The size cutoff may depend on $k$ and on the
fixed size factor.

The logical tape has read-only groups $s_i\in\{0,1\}^L$ for $i\in[k]$,
a work group $w$ indexed by $K_1$, and legacy bits $e_z$ indexed by
$z\in C$.  Each $s_i$ is embedded by the same fixed injection of $[L]$
into $K_1$; other grid entries are specified zeros.  Stages
$r=0,\ldots,M-1$ have labels $c=\gamma^r$.  Work is blank at stage zero.
During step $r=j-1\to j$, for $1\le j\le L$, copy $s_i[j]$ to a
distinct work location $w_{i,j}$ for every $i$ in parallel.  Thereafter
assign one new gate location per step, using earlier work locations or
public constants.  Reserve a fixed grid position $y_{\rm out}$ solely
for the final output assignment.  Pad remaining steps by idle work rules.
The number of work positions and nonidle steps is at most
\[
 kL+O(1)+O(L^{c_{\rm allow}}),
\]
with constants allowed to depend on the circuit size factor.  The choices
above make both counts fit below the available grid and stage sizes.

At stage $c$, require $e$ to be one-hot at $c^{-1}$.  Require every
unassigned work position to be zero, every assigned copy to equal its
specified input, and every assigned gate to satisfy its rule.  Consequently
each Boolean tuple $s=(s_i)_i$ has exactly one valid tape
$b=(w,e,(s_i)_i)$ at every stage.  Join consecutive stages by their
assignment steps, and disable the wrap step $\gamma^{-1}\to1$.

\begin{lemma}[Validity and a constant edit alphabet]\label{ia:validity}
The union of the stage-valid tapes is described by polynomially many
Boolean equations of degree at most $d=3$.  For fixed $k$, every macro
step and its inverse use a fixed finite alphabet of flips, whose
coefficients are determined by a bounded number of logical reads and
public rule bits.  Equations involving $s_i$ and other groups can be
written as the sum of an $(e,s_i)$ contribution and an $(e,w)$
contribution.  No equation requires a product involving two distinct
positive port groups.
\end{lemma}

\begin{proof}
Impose bitness on all logical coordinates, and
\begin{equation}
 e_z e_{z'}=0\quad(z\ne z'),\qquad
 \sum_{z\in C}e_z=1.                              \label{ia:one-hot}
\end{equation}
For each stage $c$, multiply all its work requirements by
$e_{c^{-1}}$.  Gate equations have degree at most two before this
multiplication, so the resulting degree is at most three.  A copy
requirement at an occupied position is specifically
\begin{equation}
 e_{c^{-1}}w_{i,j}-e_{c^{-1}}s_i[j]=0.             \label{ia:copy-equation}
\end{equation}
The requirements are vacuous at unselected stages.  At the one selected
stage, the prefix equations determine all assigned work positions in
topological order and force every other position to zero.  This proves
the characterization and the dependency assertion.

For start frame $c$, list $k$ work write slots and the two legacy slots
$c^{-1},(\gamma c)^{-1}$.  During a gate step only the first work slot
is active; use a zero unit vector for every inactive slot.  Write
$v_1(c),\ldots,v_{k+2}(c)$ for these public vectors in the logical tape
space, and for $\sigma\in\{0,1\}^{k+2}$ put
\begin{equation}
 h_\sigma(c)=\sum_{r=1}^{k+2}\sigma_r v_r(c).
                                                               \label{ia:flip-mask}
\end{equation}
At a valid pre tape, choose a work coefficient to be the old bit plus
its required new assignment; choose both legacy coefficients to be one.
The required assignment is a read input, a degree-at-most-two gate of
earlier work bits, or the unchanged value for an inactive slot.  At a
valid post tape, the inverse clears precisely the newly assigned
positions and flips the same legacy positions.  Earlier operands have
not changed, so the two operations are reciprocal.  There is a bounded
list of old-value, operand, and input reads for fixed $k$.  Public phase,
gate, activity, and wrap bits select their roles.  Define the same
coefficient rules on arbitrary Boolean read answers; at the disabled
wrap use an idle work proposal.  This gives a finite branch system
without any assumption that those arbitrary answers describe a valid
tape.
\end{proof}

\begin{definition}[Extended inner graph]\label{ia:extended-path}
An extended label is $(c,b,\eta,t)$, where $c\in C$,
$\eta\in\{0,1\}$, and $0\le t\le q-1$.  In pre mode $\eta=0$,
$b$ must be valid at stage $c$; in post mode $\eta=1$, it must be valid
at stage $\gamma c$.  Edges are the following:
\begin{enumerate}
\item in pre mode, increase $t$ by one until $q-1$;
\item at full pre progress, perform the enabled macro edit and change
to post mode, retaining the frame and full progress;
\item in post mode, decrease $t$ by one until zero;
\item at empty post progress, change to empty pre mode in frame
$\gamma c$, leaving $b$ unchanged.
\end{enumerate}
The fourth operation is included also for $c=\gamma^{-1}$; only the
second operation is disabled at that frame.
\end{definition}

\begin{lemma}[Endpoints and output]\label{ia:path-endpoints}
For each $s$, the extended labels with that read-only tuple form a
single directed path, without isolated vertices or self edges.  Its
unique source and immediate successor are respectively
\begin{equation}
 (\gamma^{-1},b_0(s),1,q-1),\qquad
 (\gamma^{-1},b_0(s),1,q-2),                      \label{ia:source-successor}
\end{equation}
where $b_0(s)$ has blank work and legacy bit $e_1=1$.  Its other
endpoint is the full pre label at stage $M-1$, and there
$w[y_{\rm out}]=f(s)$.
\end{lemma}

\begin{proof}
Replace the unique macro tape at stage $x$ by its post labels in frame
$\gamma^{-1}x$, in decreasing order of $t$, followed by its pre labels
in frame $x$, in increasing order of $t$.  The frame-advance edge joins
the two halves.  Each enabled macro edit joins the full pre end of one
such segment to the full post start of the next.  The only absent join
is the wrap join from the final stage to stage zero.  This gives exactly
the stated path and endpoints.  Every segment has more than one vertex,
so there are no isolated vertices; each edge changes a label component,
so there are no self edges.  The final work assignment and all its
ancestors have been completed at the last stage, proving the output
claim.
\end{proof}

\subsection{Chart conventions and lifted cyclic interpolation}
\label{ia:lifts-and-masks}

A native chart is an $\mathbb F_q$-coefficient polynomial on
$O(\ell)$ coordinates over $\mathbb F_q$, possibly with slots.  A
persistent chart has at most one slot axis of length $O(\ell)$; stacks
will have a bounded number of such axes.  A full rotating or projection
variable always occupies one whole $K$ label.  Native charts will later
be embedded in one common $K^u$, with unused coordinates ignored.

For slot lengths $t_1,\ldots,t_r$ and masks
$L_b\in\mathbb F_q^{t_b}$, write
\begin{equation}
 L\cdot F(Y)=\sum_{z_1,\ldots,z_r}
             \left(\prod_{b=1}^r L_b[z_b]\right)
             F_{z_1,\ldots,z_r}(Y).                \label{ia:contraction}
\end{equation}
There are $q^{O(\ell)}$ mask tuples.  To query any specified mask tuple
smoothly, choose independent uniform $A_b$, substitute
$L_b=A_b+(L_b-A_b)$, and expand the $2^r$ terms.  Each term has uniform
mask marginals, even conditional on the point and the requested tuple.
A nonzero tensor remains nonzero after independent random contractions
with probability at least $(1-1/q)^r$: at each step a nonzero linear
map has a kernel of relative size at most $1/q$.

We repeatedly use the elementary polynomial zero bound: a nonzero
polynomial of total degree $d_*$ on $\mathbb F_q^m$ vanishes on at most
a $d_*/q$ fraction of the grid.  For completeness, expand in the last
variable with leading coefficient of degree at most $d_*-r$ and last
variable degree $r$.  Inductively that coefficient vanishes on at most
$(d_*-r)/q$ of the earlier coordinates; elsewhere there are at most
$r$ roots in the last variable.  Induction on $m$ proves the bound.
We apply it also after adjoining the finitely many mask groups.
Thus any nonzero degree-$O(D+\ell h)$ chart identity, including a
whole slot list, is detected with constant probability.  An ignored
coordinate group is checked by comparing two independent replacements
of that group with common uniform other coordinates.  Its difference
is a polynomial; the same bound gives an exact alternative between
independence and constant rejection probability.  Throughout, smoothness
means bounded marginal density relative to uniform point and mask
coordinates.  Ignored coordinates are independently filled uniformly.

Ordinary interpolation on the $H$ digits extends a function on $K_1$
with individual degrees below $h$ and total degree at most
$\ell(h-1)$.  The punctured cyclic grid requires a different extension.

\begin{lemma}[Cyclic lift, including descent at zero]\label{ia:cyclic-lift}
For $0\le j<M$, write $j=\sum_{t=0}^{\ell-1}j_t h^t$ and define
\begin{equation}
 \widetilde j=\sum_{b=0}^{\ell-1}j_{vb\bmod\ell}q^b.
                                                               \label{ia:lift-exponent}
\end{equation}
Replacing $Z^j$ by $Z^{\widetilde j}$ in interpolation on $C$ gives a
linear extension to $K$, of coordinate degree at most $\ell(h-1)$.
For $\mathbb F_q$-valued data this extension has coefficients in
$\mathbb F_q$ in the chosen coordinates, including its value at zero.
It commutes with $Z\mapsto xZ$ for every $x\in C$.
\end{lemma}

\begin{proof}
Modulo $M$ we have $q^b=h^{vb}$ and $h^\ell=1$, hence
\[
 \widetilde j\equiv
 \sum_b j_{vb\bmod\ell}h^{vb\bmod\ell}=j\pmod M.
\]
The two monomials therefore agree on $C$, where the powers
$1,Z,\ldots,Z^{M-1}$ form the usual interpolation basis.  Each map
$Z\mapsto Z^{q^b}$ is linear over $\mathbb F_q$, so the coordinate
degree of the lifted monomial is at most $\sum_t j_t\le\ell(h-1)$.

Multiplication of $\widetilde j$ by $q$ modulo $q^\ell-1$ cyclically
rotates its base-$q$ digits.  Under \eqref{ia:lift-exponent}, this is
the permutation $j\mapsto qj\bmod M$.  The digit list all of whose
entries equal $h-1$ is excluded by $j<M$; exponent zero is consequently
the unique representative of residue zero on $C$.  If the original
interpolant is $\sum_j a_jZ^j$ and its values belong to $\mathbb F_q$,
uniqueness of interpolation on $C$ gives
\[
 a_{qj\bmod M}=a_j^q.
\]
The digit rotation then proves that its lift $F$ satisfies
$F(z)^q=F(z)$ for every $z\in K^\times$.  At zero, $F(0)=a_0$ and
$a_0^q=a_0$, so the same assertion holds.

The coordinate degree is below $q$ after our cutoff.  Conjugating the
coordinate coefficients by the $q$-power map leaves the values on
$\mathbb F_q^\ell$ unchanged.  Uniqueness of interpolation with
individual degrees below $q$ forces each coefficient to be fixed by
that conjugation, and thus to lie in $\mathbb F_q$.  Finally
$x^{\widetilde j}=x^j$ for $x\in C$, which proves the scaling assertion
on every basis monomial.  These arguments apply coefficientwise in any
additional independent variables.
\end{proof}

\subsection{Canonical certificates for linear conditions}
\label{ia:linear-certificates}

Partial-sum polynomial certificates also appear in Rubinstein's
construction~\cite[Sections~6.2--6.3]{Rubinstein2016Settling}.
Here each input chart has its own canonical certificates, and contributions
from different charts meet only in a tested identity.  We require these
recipes to remain linear in their input chart and, in rotating calls, to
respect cyclic scaling of the designated variables even when the final
condition fails.

\begin{lemma}[Smooth linear verification]\label{ia:linear-tool}
Let a fixed finite list of primary native charts, each scalar or with
one $O(\ell)$ slot axis, be subject to an explicit
$\mathbb F_q$-linear condition $\mathcal S$ on their full evaluation
arrays.  Its matrix has polynomially many rows and columns, with
logarithms $O(\ell\log q)$.  Coefficient and special-grid conditions
may be expressed by public linear matrices.  A nonrotating call may
adjoin a separately enforced primary $1$ for affine terms.

In a rotating call, at most one full variable in each primary is
designated to scale by a common $x\in C$.  Assume $\mathcal S$ is
invariant under this simultaneous pullback and imposes the lifted
shape of Lemma~\ref{ia:cyclic-lift} in every rotating variable.  No
arbitrary noninvariant affine term is allowed in such a call.

The condition has a canonical verification construction with the
following properties.
\begin{enumerate}
\item A fixed number of auxiliary families is stored \emph{separately
for each primary block}.  Each has native dimension $O(\ell)$ and at
most one slot axis of length $O(\ell)$.  Its canonical arrays are
linear in that primary's arrays, also when $\mathcal S$ fails.
\item Honest input coordinate degree $d_{\rm in}$ produces auxiliary
degree $O(d_{\rm in}+\ell h)$.  A fixed mixture of bounded, smooth
linear field probes is perfectly complete on legal canonical tuples
within the degree budgets.  On genuine degree-at-most-$D$ polynomial
tuples, either $\mathcal S$ holds and every auxiliary is canonical, or
the mixture rejects with a fixed positive probability.
\item In a rotating call each auxiliary has a new designated full
variable $R$, and rotating its primary by $x$ changes the auxiliary by
$R\mapsto xR$.  Its canonical recipe has the lifted shape in $R$
whenever its primary has the required shape, even if final rows fail.
All canonical arrays and public maps are polynomial-time constructible.
\end{enumerate}
\end{lemma}

\begin{proof}
We give the construction, uniqueness, degree, and probe arguments.

\paragraph{Why arbitrary scaling factors preserve the linear condition.}
In a rotating call substitute $RZ$ for the designated variable of a
primary and denote the resulting chart by $f_R$.  An unchanged primary
stays unchanged.  We first show that legal input arrays satisfy
$\mathcal S$ for every $R\in K$, rather than only for $R\in C$.
Extend scalars in the evaluation space and its linear subspace
$\mathcal S$ from $\mathbb F_q$ to $K$.  Since $|C|=M$ is odd, it is
invertible in characteristic two.  The character projections
\[
 \Pi_j=\frac1M\sum_{x\in C}x^{-j}\rho(x)
 \qquad(0\le j<M)
\]
for the simultaneous scaling action $\rho$ decompose the space and
preserve every invariant subspace.  On the imposed lift shape, the
character $j$ is represented by the single exponent $\widetilde j$
in every rotating primary.  Nonrotating components occur in character
zero.  Substitution by $R\ne0$ therefore acts by the scalar
$R^{\widetilde j}$ on each character space, and preserves the extended
subspace.  At $R=0$ it projects onto character zero, which also preserves
the subspace.  Scaled evaluations of our $\mathbb F_q$-valued chart
are still $\mathbb F_q$-valued; intersecting with the original
evaluation space, or using the original defining rows, proves the
assertion over $\mathbb F_q$.

A failed original row fails at $R=1$, so it remains a nonzero polynomial
condition.  Tests requiring the substitution sample $R$ uniformly from
$K^\times$; this changes a uniform density by at most
$(1-|K|^{-1})^{-1}$ and keeps primary point marginals uniform.  The
substitution is bilinear in coordinate variables and at most doubles
total degree.  In a nonrotating call omit $R$ everywhere below.

\paragraph{Separate row contributions and partial sums.}
Pad the row set by zero rows to $H^\rho$, with $\rho=O(\ell)$.
For one primary block with $m_0$ native $q$-coordinates, put $m'=vm_0$.
An $H$-basis of $\mathbb F_q$ gives a packing map
\[
 w_{\rm pack}:\mathbb F_q^{m'}\longrightarrow\mathbb F_q^{m_0}.
\]
This is a surjective $\mathbb F_q$-linear map whose restriction to
$H^{m'}$ is a bijection onto $\mathbb F_q^{m_0}$.  In particular it
sends uniform inputs to uniform outputs.  For each primary slot $s_0$,
interpolate that block's public matrix entries as
$A_{s_0}(I,X)$ on $H^\rho\times H^{m'}$, with individual degrees below
$h$.  Define, for this block only,
\begin{align}
 p_{m'}(R,I,X)
   &=\sum_{s_0}A_{s_0}(I,X)
                 f_{s_0,R}(w_{\rm pack}(X)),       \label{ia:linear-top}\\
 p_j(R,I,X_1,\ldots,X_j)
   &=\sum_{x_{j+1},\ldots,x_{m'}\in H}
      p_{m'}(R,I,X_1,\ldots,X_j,x_{j+1},\ldots,x_{m'})
      \quad(0\le j<m').                           \label{ia:partial-sums}
\end{align}
The top identity requests the primary with the public slot mask
$(A_{s_0}(I,X))_{s_0}$, made smooth by additive splitting.  This is a
linear probe; the matrix entries are not unknown operands.

The final condition is the tested identity
\begin{equation}
 \sum_{\text{primary blocks }b} p_{0,b}(R,I)=0.
                                                               \label{ia:final-row-sum}
\end{equation}
\emph{No chart for this mixed sum is stored.}  On legal data it
vanishes on every $I\in H^\rho$ and every $R\in K$.  Its individual
$I$ degrees are below $h$, and its honest $R$ coordinate degrees are
below $q$, so this is a polynomial identity.  Conversely, its
vanishing at $R=1$ and at the row grid, with correct partial sums,
is exactly the original matrix condition.

\paragraph{Division enforces all partial sums uniquely.}
For $1\le j\le m'$ introduce quotient and remainder charts and impose
\begin{align}
 p_j&=(X_j^h-X_j)Q_j+X_j^{h-1}p_{j-1}+E_j,
                                                        \label{ia:division}\\
 p_{j-1}&\text{ is independent of }X_j,
 \qquad \deg_{X_j}E_j\le h-2.                       \label{ia:division-bounds}
\end{align}
In characteristic two,
\[
 \sum_{x\in H}x^r=0\quad(0\le r<h-1),\qquad
 \sum_{x\in H}x^{h-1}=1.
\]
Here the $r=0$ identity is $h=0$ in the field.  The coefficient of
$X_j^{h-1}$ in the remainder modulo $X_j^h-X_j$ is thus precisely the
sum of the original polynomial over $X_j\in H$.  Euclidean division
by this monic polynomial uniquely fixes $Q_j$, the top remainder
coefficient $p_{j-1}$, and $E_j$.  Starting at
\eqref{ia:linear-top}, downward induction fixes all certificates and
gives \eqref{ia:partial-sums}.  Independence from later trailing
coordinates follows in the same induction.

To test the degree restriction without concentrated reads, write
\begin{equation}
 G_j(R,I,X_{\ne j},U)
   =\sum_{b=0}^{h-2}E_{jb}(R,I,X_{\ne j})U_b,
 \qquad U\in\mathbb F_q^{h-1},                     \label{ia:coefficient-chart}
\end{equation}
where $E_{jb}$ is the coefficient of $X_j^b$ in $E_j$.  Check
independence from $X_j$, addition in $U$, and homogeneity under a
uniform nonzero scalar $\lambda\in\mathbb F_q$.  Sampling a nonzero
scalar makes $\lambda U$ uniform; zero homogeneity already follows
from addition.  If these identities are exact, $G_j$ is pointwise
$\mathbb F_q$-linear in $U$, and uniqueness of reduced interpolation
makes its polynomial linear in $U$.  Now check
\begin{equation}
 E_j(R,I,X)=G_j(R,I,X_{\ne j},(1,X_j,\ldots,X_j^{h-2})).
                                                               \label{ia:coefficient-evaluation}
\end{equation}
The special $U$ tuple is queried as two uniform additive summands.
Only \emph{after} linearity has been established in the exact-soundness
case analysis do we use the $O(D+h)$ degree bound for this substitution.
The earlier linearity discrepancies have degree $O(D)$ and themselves
have constant gap.  Equation~\eqref{ia:coefficient-evaluation} then
enforces \eqref{ia:division-bounds} and fixes $G_j$ uniquely.

\paragraph{A whole list is checked at once.}
Sampling a single recurrence index would lose a factor $m'$.  Instead
choose a cyclic permutation $\omega$ of the $m'$ scalar coordinates
so that coordinate $j$ of $\omega^jX$ is one fixed coordinate of $X$.
Store
\begin{equation}
 \overline p_j(R,I,X)=p_j(R,I,\omega^jX),
 \qquad
 p_{j-1}(R,I,\omega^jX)
      =\overline p_{j-1}(R,I,\omega X),             \label{ia:recurrence-frames}
\end{equation}
and store $Q_j,E_j,G_j$ in the equation-$j$ frame, with $U$ unchanged.
Both public powers in \eqref{ia:division}, and the predecessor
independence axis in its corresponding frame, are now common across
all $j$.  A uniform mask on the recurrence-index list checks all its
equations with boundedly many point probes.  Endpoint slices and
predecessor shifts only change the requested linear mask, which is
split smoothly.  The coefficient families use the same $U$ group in
every slot.  If $m'=0$, retain only the unsummed top contribution and
introduce no division family.

\paragraph{Degree, covariance, and exact gap.}
The top degree is at most
\begin{equation}
 2d_{\rm in}+(\rho+m')(h-1).                       \label{ia:tool-degree}
\end{equation}
Monic division by $X_j^h-X_j$, reduction, coefficient extraction, and
finite summation do not increase total degree.  Multiplication by a
linear $U_b$ costs at most one.  This proves the honest degree bound.
There is a fixed number of families per primary; the growing recurrence
index is one $O(\ell)$ slot axis, and the additional $U$ domain has
$h-1\le\ell$ coordinates.

First test ignored coordinates and the independence and linearity
devices.  Unless they hold identically, one has a nonzero polynomial
discrepancy of degree $O(D)$ and constant rejection probability.
Once they hold, every remaining displayed discrepancy has degree
$O(D+\ell h)=o(q)$ and a failed slot identity survives a random
contraction with constant probability.  Exact identities give the
uniqueness just proved and then \eqref{ia:final-row-sum}.  A fixed
mixture of test types therefore has a uniform positive gap.  Perfect
completeness concerns canonical tuples within the budgets; it does not
claim that every legal primary of degree exactly $D$ has auxiliaries
of degree at most $D$.

For covariance even on invalid inputs, expand a lifted primary after
substitution over $K$ as
\[
 f_R=\sum_{j=0}^{M-1}R^{\widetilde j}v_j
\]
in its evaluation arrays.  All subsequent operations in
\eqref{ia:linear-top}--\eqref{ia:coefficient-chart} are the same fixed
linear maps on that primary, independent of $R$.  This includes
division and coefficient extraction, since honest-degree evaluation
arrays determine the required polynomials by interpolation.  These
maps preserve the exponent list, and rotating the input by $x\in C$
replaces $R$ by $xR$.  No final zero-row identity is used.  In particular
every canonical certificate is defined even when its primary belongs
to data failing the final condition.

Finally, all point substitutions used in these tests are coordinate
permutations, surjective packing, nonzero scaling, fresh independent
coordinate replacement, or additive splitting.  They have bounded
point marginals.  Requested masks are split.  All arrays have
polynomial size, and every construction operation is finite linear
algebra on them.  This proves the remaining claims.
\end{proof}
\subsection{Persistent charts and their dependency tags}
\label{ia:persistent}

Every persistent chart has a canonical recipe $F(c,b;Y)$, possibly
with one slot axis.  At each fixed frame $c$, both its values and its
coordinate coefficients have an assigned multilinear Boolean degree
$\kappa_F\le d=3$ in $b$, after reducing by $b_j^2=b_j$.  This
logical degree is different from the coordinate degree bounded by
$D$.  All recipes are defined for every Boolean tape, including tapes
that fail validity.  A recipe either is independent of $c$, or has
one designated full variable whose scaling realizes $c\mapsto xc$.

Give each family a tag in $\{0,1,\ldots,k\}$.  A tag-$0$ recipe
depends only on $c,w,e$ and public data; that public data may include
the circuit, work addresses, and the port index.  A tag-$i$ recipe
depends only on $c,s_i,e$ and public data fixed by the standardized
one-port template.  It may ignore $s_i$ if it is a certificate created
inside that template.  Call tags compatible if they are equal or one
of them is zero.  After fixing a finite Boolean transition branch,
tag-$0$ update operands will all have tag $0$, and tag-$i$ update
operands will have compatible tags.

\paragraph{Monomials and the separated validity calls.}
Use the ordinary full-grid extension for $W_w$ and each $W_{s_i}$,
and the cyclic lift for $W_e$.  Include the constant chart $1$,
tested against its prescribed value.  For all ordered lists of length
at most $d$ from $\{w,e\}$, and separately from $\{s_i,e\}$ for
each $i$, include
\begin{equation}
 W_{g_1,\ldots,g_j}(b;Y_1,\ldots,Y_j)
    =\prod_{r=1}^j W_{g_r}(b;Y_r).                 \label{ia:monomial-charts}
\end{equation}
Repeated types use independent variables.  Share the legacy-only
types.  A monomial containing $s_i$ has tag $i$; all other monomials
have tag $0$.  Enforce each nontrivial product by factorization into
its shorter prefix and last factor on independent native variables,
together with all ignored-coordinate requirements.

Apply the nonrotating linear tool as follows.
\begin{enumerate}
\item One tag-$0$ validity call uses the monomials from $\{w,e\}$
and $1$.  It imposes the order-one extension shapes, bitness on the
grids, the one-hot requirements \eqref{ia:one-hot}, and all blank and
gate requirements of Lemma~\ref{ia:validity}.  Use one work-requirement
row for each pair (stage, work grid position), with a zero row when
that position is an occupied copy.  Other conditions have fixed
lists of rows.  Circuit data affect matrix entries only.
\item For each $i$, a separate copy/port call imposes the shape,
padding, and bitness of $W_{s_i}$, and the copy rows indexed by
$(c,j)$ for $c\in C$ and $1\le j\le L$.  The row is zero when
the stage index of $c$ is less than $j$.  Otherwise it compares
\begin{equation}
 W_{e,s_i}(c^{-1},\iota(j))
     -W_{e,w}(c^{-1},w_{i,j})=0,                  \label{ia:separate-copy-row}
\end{equation}
where $\iota$ is the common input injection.  Primary-list, row-order,
and padding conventions are identical in every such call.  The
matrix blocks for port-dependent primaries depend only on these
common stage and source-bit indices.  Work destinations occur only
in the separate $W_{e,w}$ matrix block.  Auxiliaries for port-dependent
blocks receive tag $i$; auxiliaries for the work-side block receive
tag $0$.
\end{enumerate}
For example, bitness is the linear row
\[
 W_{g,g}(y,y)-W_g(y)=0
\]
on represented grid entries.  Although its two grid arguments
coincide, it is verified as a matrix row, not by a concentrated
unknown-product query.  The same representation handles every
degree-at-most-three Boolean validity term.  Extension shapes are
linear evaluation/interpolation conditions; the cyclic shape is
$\mathbb F_q$-linear by Lemma~\ref{ia:cyclic-lift}.

If these conditions and the factorization identities hold exactly
on genuine polynomials, order-one grid entries are bits, their
shapes fix their full extensions, products fix higher monomials,
and the rows give exactly a valid union-of-stages tape.  The linear
tool fixes each certificate separately.  The initial auxiliary
assigned degrees can all be bounded by $d$, or by the smaller
degree of their primary when useful.  Since every certificate map
is linear in one primary, the same degree statement holds on
arbitrary Boolean tapes even when a final validity row is nonzero.

Include also a tag-$0$, point-constant output chart
\begin{equation}
 F_{\rm out}(b)=w[y_{\rm out}].                    \label{ia:output-chart}
\end{equation}
A separate nonrotating linear call ties it to the designated grid
entry of $W_w$ and enforces its ignored coordinates.  Give it and
the associated certificates assigned degree one.  This output chart
does not need a critical-address projection for transitions.

\paragraph{Frame and clock-times-monomial charts.}
Let $J_c(X)$ be the cyclic lift of the indicator of $c^{-1}$ on
$C$.  It has tag $0$, assigned degree zero, and
\begin{equation}
 J_{xc}(X)=J_c(xX).                                \label{ia:clock-covariance}
\end{equation}
Tie it to the legacy chart by
\begin{equation}
 W_e(X)=J_c(\gamma^\eta X).                        \label{ia:clock-legacy}
\end{equation}
Below, $\eta$ is enforced as a point-constant bit.  Equation
\eqref{ia:clock-legacy} is then a finite choice of two smooth
polynomial identities.  One-hot legacy validity determines a unique
$c$ and the required pre/post relation, and fixes $J_c$.  For every
monomial $W$, including $1$, include
\begin{equation}
 T_W(c,b;X,Y)=J_c(X)W(b;Y),                         \label{ia:clock-monomial}
\end{equation}
with independent arguments.  It has the same tag and degree tier as
$W$, and rotates only $X$.  Enforce it by factorization.

\begin{lemma}[Clock-indexed linear formulas]\label{ia:clock-formula}
Suppose for each $y\in C$ that $L_y$ is an explicitly given
$\mathbb F_q$-linear map from selected monomial evaluation arrays
of logical degree at most $k_0$ to the full native array of a scalar
or one-slot output.  There is a canonical chart, lifted in $Z$, with
\begin{equation}
 G(c,b;z,\cdot)=L_{cz}(W(b))\quad(z\in C).         \label{ia:clock-formula-value}
\end{equation}
It rotates only $Z$, and it and its verification auxiliaries have
assigned degree at most $k_0$.  A rotating linear-tool call verifies
them.  The recipes and their covariance hold for all Boolean $b$.
\end{lemma}

\begin{proof}
Use the needed $T_W$ inputs and the output $G$ as primaries.  Impose
their lifted shapes in the designated variables and the linear rows
\begin{equation}
 G(z,\cdot)=\sum_{x\in C}L_{x^{-1}z}(T(x,\cdot))
 \quad(z\in C).                                  \label{ia:clock-convolution}
\end{equation}
Rows also range over every output evaluation position and slot.
On canonical $T$ inputs, only $x=c^{-1}$ contributes, giving
\eqref{ia:clock-formula-value}.  The lift shape then determines the
whole chart uniquely.  The row condition is invariant under the
simultaneous replacement $G(z)\mapsto G(u_0z)$ and
$T(x)\mapsto T(u_0x)$: setting $x'=u_0x$ changes the map index to
$(x')^{-1}u_0z$.  The shape conditions are invariant too, so
Lemma~\ref{ia:linear-tool} applies.  Every input has logical degree
at most $k_0$; interpolation and all per-primary certificate operations
are linear and preserve that degree.  Covariance follows first on
$C$ and then by uniqueness of the lift, and for the auxiliaries by
the last part of the linear tool.  No validity row on $b$ was used.
Public clock-dependent terms are supplied through $T_1=J_c$, keeping
the rotating condition linear.
\end{proof}

Apply this lemma to a fixed list, for fixed $k$, of critical addresses
needed by either transition direction.  A critical output chart has
only the variable $Z$ and gives at $z\in C$ the required logical read
or rule bit for absolute frame $cz$.  For each port, make a separate
tag-$i$ read chart that reads $s_i[j]$ during initial copy step $j$
and is identically zero on every other step.  This is a standardized
degree-one call independent of all subsequent computation.  Work
reads have separate tag-$0$ charts, defined as zero when inactive;
there are a bounded number per write slot for old and operand reads.
Phase, gate, activity, and wrap-disable bits use tag-$0$ degree-zero
calls.  Fix their numbers and tiers by the rule basis and write-slot
bound, padding inactive roles.  Their public values and work address
maps may depend on the circuit.

Every auxiliary of a tag-$0$ clock call has tag $0$.  Every auxiliary
of a standardized tag-$i$ clock call has tag $i$, even an auxiliary
for a shared legacy-only or public $T$ primary in that call.  The
shared primary itself retains its base tag and formula.  This
distinction between an input's tag and the ancestry of its
per-call certificate is part of the construction.

A fixed-frame tape edit must update every persistent chart, including the
certificates introduced to verify other charts.  Its effect on a chart
$F$ is the difference
\[
 F(c,b+h_\sigma(c);Y)-F(c,b;Y).
\]
For a positive-tier chart we store this difference in a new chart with a
clock variable $Z$, whose restriction at $Z=1$ gives the required change.
Taking a Boolean difference lowers the logical degree.  We can also store the
changes of this new chart and its certificates, stopping after at most
three levels.  The construction below supplies these polynomial charts;
the projection stacks in Section~\ref{ia:projection-stacks} will realize
the restriction at $Z=1$ without a concentrated point read.

\begin{lemma}[Derivative closure for all persistent families]
\label{ia:derivative-closure}
The persistent system can be closed under every required fixed-frame
Boolean edit, including edits to verification auxiliaries, with a
fixed number of families for fixed $k$.  Each has at most one
$O(\ell)$ slot axis, native dimension $O(\ell)$, honest coordinate
degree $O(\ell h)$, and assigned logical degree at most three.  The
closure and its polynomial-time canonical recipes are defined on
invalid Boolean tapes as well.
\end{lemma}

\begin{proof}
Let $F$ have positive assigned tier $k_0$.  For every required flip
mask, retain its native variables $Y$ and slots, adjoin a new full
variable $Z$, and prescribe on the cyclic grid
\begin{equation}
 G_F^\sigma(c,b;z,Y)
   =F(cz,b+h_\sigma(cz);Y)-F(cz,b;Y),\qquad z\in C.
                                                               \label{ia:derivative}
\end{equation}
Extend in $Z$ by the cyclic lift.  Both parent terms are evaluated
in frame $cz$; the parent's existing argument coordinates remain
exactly $Y$.  A tag-$0$ parent uses the full flip alphabet
\eqref{ia:flip-mask}.  For a tag-$i$ parent, only the two legacy
coefficients matter: the input bits $s_i$ are never written and
work bits are absent from its recipe.  Use the fixed four masks
\begin{equation}
 h_\chi(y)=\chi_1\mathbf 1_{e_{y^{-1}}}
            +\chi_2\mathbf 1_{e_{(\gamma y)^{-1}}},
 \qquad \chi\in\{0,1\}^2,                        \label{ia:port-flip-mask}
\end{equation}
independent of $k$ and $i$.

For a squarefree monomial, a fixed Boolean translation gives
\begin{equation}
 \prod_{j\in S}(b_j+h_j)-\prod_{j\in S}b_j
   =\sum_{T\subsetneq S}\left(\prod_{j\in T}b_j\right)
                         \left(\prod_{j\in S\setminus T}h_j\right).
                                                               \label{ia:degree-drop}
\end{equation}
Thus \eqref{ia:derivative} has logical degree at most $k_0-1$.
Lifting its new variable adds at most $\ell(h-1)$ to its coordinate
degree.  For each absolute frame $y$, output evaluation point, and
inherited slot, the derivative is a known linear combination of
grid monomials of degree at most $k_0-1$: from $\{w,e\}$ for tag
$0$, or from $\{s_i,e\}$ for tag $i$.

These coefficients are computed from the \emph{canonical parent
recipe}, not from the presently supplied parent table.  An explicit
method is bounded-degree Boolean interpolation.  For a multilinear
function $P$ of degree at most $r$ on a universe of $N$ bits, its
coefficient on $S$, $|S|\le r$, is
\begin{equation}
 [b_S]P=\sum_{T\subseteq S}(-1)^{|S|-|T|}P(\mathbf 1_T).
                                                               \label{ia:boolean-inversion}
\end{equation}
In characteristic two the signs are all one.  Use this on the
right side of \eqref{ia:derivative}, which is computable on every
Boolean argument, including translated invalid tapes.  There are
only $\sum_{j\le r}\binom Nj=\poly(n')$ required subsets, because
$r\le3$ and $N=\poly(n')$.  For tag $i$ the interpolation universe
is the common ordered list of its $L$ source bits and $M$ legacy
bits, with no work coordinates.  Interpret the resulting coefficients
as public maps $L_y$ on the permitted monomial arrays.

Apply Lemma~\ref{ia:clock-formula} using only inputs through tier
$k_0-1$.  The output derivative and every newly introduced auxiliary
receive that tier and the parent's tag.  In particular a certificate
for a shared tag-$0$ $T$ input in this positive call receives the
positive ancestry tag.  The derivative rotates only its new variable,
since on the grid
\[
 G_F^\sigma(u_0c,b;z,Y)=G_F^\sigma(c,b;u_0z,Y),
\]
and lift uniqueness extends this to $K$.  Its older parent variables
are not additionally scaled under a frame change.  Each verification
auxiliary instead rotates its own new $R$, as proved by the linear
tool.

Repeat this operation for every added positive-tier family, including
derivatives and all their proof families.  Every descendant call
strictly lowers its assigned tier, whether or not its actual degree
was already smaller.  The recursion ends after at most three levels.
A family has only finitely many children for fixed $k$.  A derivative
inherits its parent's single slot axis; a verification auxiliary uses
its own recurrence-index axis rather than tensoring it with primary
slots.  Consequently the number of families is fixed, their native
dimensions and slot lengths remain $O(\ell)$, and their coordinate
degrees remain $O(\ell h)$.  Tier-zero charts require no derivative
for a fixed-frame Boolean edit.

All maps have polynomial construction cost: there are polynomially
many frames, evaluations, slots, monomials, and rows, and only a
fixed number of recursive levels.  Finite sums, monic division,
coefficient extraction, and interpolation are fixed linear operations
on each primary array.  They neither create a further logical product
nor import another primary's dependency.  This remains true when the
final row sum fails, by Lemma~\ref{ia:linear-tool}.  The asserted
invalid-tape recipes and their degree bounds follow as well.
\end{proof}

\begin{lemma}[Literal positive-port ancestry]\label{ia:port-ancestry}
There are deterministic local call and ancestry conventions such that
the entire tag-$i$ construction in an actual $k$-port graph is
identical, after renaming $s_i$ to $s$, to the positive-tag construction
in the one-port template with the same $L,n'$ and field parameters.
The agreement includes primary lists, row indexing and zero padding,
packing maps, all auxiliary and derivative families, native dimensions,
slot orders, and canonical arrays whenever $(c,e,s_i)=(c,e,s)$.
It holds independently of the actual work destinations, $k$, and $f$.
\end{lemma}

\begin{proof}
Choose row lists, their order, and zero padding separately for each
call's local template.  Do not choose a row dimension by maximizing
over all calls in the actual graph.  In each copy/port call use the
same input-shape rows, bitness rows, and $(c,j)$ copy rows, and the
same full primary-list conventions.  Although the work-side matrix
depends on $w_{i,j}$, its dimensions do not.  This makes all
port-dependent primary blocks literal copies of their dummy blocks.

For a concrete example, consider $W_{e,s_i}$ in
\eqref{ia:separate-copy-row}.  Its matrix $A_{e,s}(I,X)$ knows the
stage and the source position $\iota(j)$, and has no reference to
$w_{i,j}$.  That destination appears only in the separate
$A_{e,w}$ block.  The top formula for its positive certificate is
\[
 p_{m',e,s}(I,X)
       =A_{e,s}(I,X)W_{e,s_i}(w_{\rm pack}(X)),
\]
with the evident slot contraction when needed.  Its row padding,
packing, recurrence permutation, recurrence-index order, and
coefficient-vector domain are therefore exactly those of the
one-port construction.  Each partial sum, quotient, remainder, and
coefficient chart uses that block alone.  The equality with the
work-side contribution is tested only in \eqref{ia:final-row-sum};
there is no mixed stored chart from which work addresses could
enter this positive family.

The critical source-bit calls agree directly, since they use the
same initial-copy address rule and the same zero rule thereafter.
For the remaining families, induct on positive ancestry.  If a
parent's recipe and native layout agree, its derivative uses exactly
the same restricted mask \eqref{ia:port-flip-mask}.  Its values on
every subset in \eqref{ia:boolean-inversion}, including translated
invalid inputs, are identical in the two templates.  Hence its maps
$L_y$ are identical.  Make each derivative call separately by parent
type and mask, with the fixed allowed lower-tier monomial list, full
output-equality rows, shape rows, deterministic order, and local
padding.  Its output and every per-primary auxiliary then agree
literally by the linear-tool formulas.  This includes auxiliaries
whose primary is a shared legacy-only or constant $T$ chart: the
input formula is common and the auxiliary's own call matrix is the
standardized positive one.  These auxiliary families are themselves
parents in the induction.  Strictly decreasing tiers terminate it.

Keep descendants from distinct standardized calls separate, without
coalescing them differently with circuit-dependent tag-$0$ families.
Additional full work masks and additional tag-$0$ calls then create
separate branches and do not change any positive branch.  The
construction does not use validity to establish this correspondence.
\end{proof}

\begin{lemma}[A common dimension and inherited full labels]
\label{ia:placements}
The common point domain $K^u$ and all per-chart dimension bounds can
be fixed before $k$.  Persistent charts can be placed in it so that
every derivative preserves all its parent's labels and slot order
and adds an unused whole label for its new variable.  Corresponding
positive families have identical placements.  For each fixed $k,L$,
all actual graph family and slot layouts can be fixed independently
of the size-bounded circuit's values and wiring.
\end{lemma}

\begin{proof}
Initially a monomial uses at most $d$ full variables and a $T$ chart
one more.  Base linear calls have native dimension $O_d(\ell)$.
Their row counts are bounded by $q^{C\ell}$ for an absolute fixed
$C$ independent of $k$: the work grid, clock grid, input length, and
stage count are at most $h^\ell$, and shape conditions are polynomial
in the corresponding full evaluation sizes.  Occupied-copy calls
are separate, so their number does not enter an individual row
dimension.  They use fixed stage/source-index lists.

For a derivative call on a parent of $m_0$ native coordinates and
one possible slot axis, its output, shape, and clock-equality rows
have size polynomial in $q^{m_0+C\ell}$, with a fixed exponent.
Its permitted monomial input list is fixed by its tag and tier.
Different masks have separate calls.  Therefore $\rho/\ell$, the
new native dimension divided by $\ell$, and the new slot length
divided by $\ell$ are bounded by a fixed multiple of one plus the
corresponding parent bounds.  Digit unpacking contributes the fixed
factor $v$; the coefficient-vector domain contributes at most
$h-1\le\ell$ coordinates.  Iterating through the at most three
tier drops gives an absolute bound $C_{\rm dim}\ell$, independent
of the number of calls or ports.

Choose $u$ to accommodate this bound and at most $d+1$ additional
whole labels along an inherited derivative-axis chain.  Put every
full $K$ variable in a whole label, and pack remaining native
$q$-coordinates into additional labels, ignoring their padding.
An initial chart or a new verification auxiliary is a placement
root; unrelated roots may reuse labels according to their native
templates.  A derivative retains exactly its parent's placements
and chooses a deterministically first unused whole label for its
new variable.  Repeat along that chain.  No labels have to be
globally disjoint across all ports: tests between unrelated charts
assemble their native arguments separately and fill ignored
coordinates uniformly.  Enforce all ignored-coordinate conditions.
These deterministic choices depend only on local templates and
ancestry, proving positive-family agreement.

For tag-$0$ calls at fixed $k,L$, use the full uniform
(stage, work-position) validity list, fixed critical-read roles and
tiers, the full fixed flip alphabet for every parent, and the
allowed monomial, shape, and output lists fixed by domain and tier.
Only matrix entries and rule values vary with the circuit.  This
makes the entire actual layout independent of circuit wiring, as
claimed.
\end{proof}
\subsection{Projection stacks and their exact tests}\label{ia:projection-stacks}

Critical reads and tape edits require algebraic restrictions of charts
with one or more full labels set to $1$.  Reading these special points
directly would not preserve a uniform point marginal.  Instead we
express each restriction as a sum of invertible affine pullbacks and
store the affine-pulled values in slots.  The progress coordinate $t$
specifies which of these slots are active.  At full progress, summing
over the new stack axes will give the required restriction; changing
progress will use bounded slot shifts and affine pulls.

Give each critical address output its clock-formula variable as a
projection axis.  Other initial persistent charts, including all
verification auxiliaries, start with no projection axes.  A derivative
inherits its parent's axes, with their exact labels, and adds its new
outer variable.  New verification auxiliaries again start with no
axes.  Denote the resulting set by $A_F$; its size is at most $d+1$.

Fix a generator $g$ of $\mathbb F_q^\times$.  For a full label
$\alpha$, let $D_\alpha$ be pullback by
\begin{equation}
 Z_\alpha\longmapsto 1+g(Z_\alpha-1).              \label{ia:dilation}
\end{equation}
This is an invertible $K$-affine map on the common domain.

\begin{lemma}[Projection by full-label dilations]\label{ia:projection}
If the coordinate degree in label $\alpha$ is below $q-1$, then
\begin{equation}
 \sum_{r=0}^{q-2}D_\alpha^r F
       =F\big|_{Z_\alpha=1},                      \label{ia:projection-identity}
\end{equation}
where the right side ignores that label.  Dilations on distinct
full labels commute.
\end{lemma}

\begin{proof}
Center the $\mathbb F_q$ coordinates of $Z_\alpha$ at the coordinate
vector of $1\in K$.  Decompose $F$ into parts homogeneous of degree
$r$ in the centered coordinates.  Their sums under all nonzero
scalar dilations have multiplier
\[
 \sum_{j=0}^{q-2}g^{jr}
   =\begin{cases}1,&r=0,\\0,&0<r<q-1.\end{cases}
\]
For $r=0$ this is $q-1=1$ in characteristic two, and for positive
$r$ it is a geometric sum with ratio different from one.  Thus
only the constant homogeneous part remains.  Pullbacks on disjoint
labels commute by direct substitution.
\end{proof}

\paragraph{Mode and progress marker.}
Represent $\eta$ by a point-constant scalar chart.  Enforce constancy,
then membership in $\mathbb F_2\cdot1$ by random additive binary
forms annihilating $1$.  A value outside this subspace is detected
with probability $1/2$; on the subspace, a fixed form sending $1$
to $1$ reads the bit.

Let $N=B^a$.  Represent progress by a point-constant marker tensor
$M^{\rm mark}_j$, with $a$ radix axes of length $B$, indexing
$0\le j<N$.  At $t=0$ it is zero; at $t>0$ it is one-hot at $t-1$.
After constancy is enforced, test with independent uniform separable
masks $L,L'$ the identity
\begin{equation}
 (L\cdot M^{\rm mark})(L'\cdot M^{\rm mark})
                  =(LL')\cdot M^{\rm mark},       \label{ia:marker-identity}
\end{equation}
where each axis of $LL'$ is an entrywise product.  Expanding in the
separate mask variables, exactness says
$M_j^{\rm mark}M_{j'}^{\rm mark}=0$ for $j\ne j'$ and
$(M_j^{\rm mark})^2=M_j^{\rm mark}$ for each $j$.  Hence the tensor
is zero or has exactly one coefficient equal to one.  A nonzero
discrepancy is a polynomial of bounded mask degree and is detected
with constant probability.  Enforce zero support for
\begin{equation}
 q-1\le j<N.                                      \label{ia:marker-padding}
\end{equation}
This range is a union of at most $a+1$ radix rectangles: compare
digits from the most significant end, separating the first larger
digit and the boundary itself.  Masked restrictions check every
rectangle with boundedly many probes.  The remaining possibilities
are exactly $t\in\{0,\ldots,q-1\}$.  Both mode and marker have
tag $0$.

For each persistent $F$ and each nonempty $S\subseteq A_F$, include
a stack with the same tag, point placements, and inherited slots as
$F$, and $a$ additional radix slot axes for each $\alpha\in S$.
Its intended formula is
\begin{equation}
 C^S_{F,\mathbf j}(c,b,\eta,t;Y)=
 \begin{cases}
 \displaystyle\left(\prod_{\alpha\in S}
           D_\alpha^{t-1-j_\alpha}\right)F(c,b;Y),
             &j_\alpha<t\text{ for every }\alpha\in S,\\[2pt]
 0,           &\text{otherwise},
 \end{cases}
 \qquad C_F^\varnothing=F.                         \label{ia:stack-formula}
\end{equation}
Indices run through $0,\ldots,N-1$ on each integer axis.

\begin{lemma}[Canonical stacks and bounded radix probes]\label{ia:stacks}
For a legal marker, the stacks are uniquely characterized by a fixed
number of types of bounded-probe polynomial identities.  These tests
are perfectly complete and have constant exact gap on
degree-at-most-$D$ polynomials.  Their only unknown products are
marker times smaller-stack products, with independent full operand
point marginals and compatible tags.  All slot shifts, restrictions,
and boundary slices have bounded separable-mask expansions.
\end{lemma}

\begin{proof}
For each nonempty $S$, fix one $\alpha\in S$ and require
\begin{equation}
 C^S_{F,(j_\alpha,\mathbf j_-)}
  =D_\alpha C^S_{F,(j_\alpha+1,\mathbf j_-)}
      +M^{\rm mark}_{j_\alpha}
                 C^{S\setminus\{\alpha\}}_{F,\mathbf j_-},
                                                               \label{ia:stack-recurrence}
\end{equation}
with the upper boundary at $j_\alpha=N$ defined to be zero.
For fixed $\mathbf j_-$ this is triangular downward in $j_\alpha$.
If the marker is zero its only solution is zero.  If it is one-hot
at $t-1$, its solution is
\[
 C^S_{F,(j_\alpha,\mathbf j_-)}
   =\begin{cases}
     D_\alpha^{t-1-j_\alpha}
             C^{S\setminus\{\alpha\}}_{F,\mathbf j_-},
                       &j_\alpha<t,\\
     0,                 &j_\alpha\ge t.
    \end{cases}
\]
Induction on $|S|$ gives exactly \eqref{ia:stack-formula}.

Here are explicit slot formulas for the shift, including boundary
handling.  Use little-endian radix digits
$j=\sum_{s=0}^{a-1}j_sB^s$ and let
$\lambda_L(j)=\prod_{s=0}^{a-1}L_s[j_s]$.  With out-of-range
coefficients zero,
\begin{align}
 \sum_{j=0}^{N-1}\lambda_L(j)C_{j+1}
   &=\sum_{r=0}^{a-1}L^{(r,+)}\cdot C,
                                                        \label{ia:radix-plus}\\
 L_s^{(r,+)}[d]
   &=\begin{cases}
      L_s[B-1]\mathbf1_{d=0},&s<r,\\
      \mathbf1_{d\ge1}L_r[d-1],&s=r,\\
      L_s[d],&s>r,
     \end{cases}                                      \nonumber\\
 \sum_{j=0}^{N-1}\lambda_L(j)C_{j-1}
   &=\sum_{r=0}^{a-1}L^{(r,-)}\cdot C,
                                                        \label{ia:radix-minus}\\
 L_s^{(r,-)}[d]
   &=\begin{cases}
      L_s[0]\mathbf1_{d=B-1},&s<r,\\
      \mathbf1_{d\le B-2}L_r[d+1],&s=r,\\
      L_s[d],&s>r.
     \end{cases}                                      \nonumber
\end{align}
Indeed the first pullback mask on index $k>0$ is
$\lambda_L(k-1)$.  Partition $k$ by its least digit $r$ which is
nonzero; precisely the lower digits borrow.  The second pullback
is $\lambda_L(k+1)$ for $k<N-1$; partition by the least digit below
$B-1$.  Overflow is absent in both formulas.  Slicing at zero is
a single separable restriction, and interval restrictions use the
bounded rectangle decomposition used in \eqref{ia:marker-padding}.
Tensoring these expansions over the bounded number of stack axes
still gives a bounded number of separable masks.  Split every
requested mask additively as in \eqref{ia:contraction}.

Contract the entire identity \eqref{ia:stack-recurrence} over all
slots.  The last term factors into a contraction of the marker on
axis $\alpha$ and a contraction of the smaller stack on its other
axes and inherited slots.  In the exact-soundness analysis marker
constancy has already been imposed, so its full point can be sampled
independently of the smaller stack point.  This retains perfect
completeness on canonical data.  If the recurrence is not an
identity, its discrepancy has point degree at most $2D$ and bounded
mask degree, giving constant rejection probability.  Honest stacks
are affine pullbacks of persistent charts or zero, so they have the
same honest coordinate degree.  Marker tag $0$ is compatible with
every parent tag.  This proves all claims.
\end{proof}

\subsection{Exact validation and locally encoded evolution}
\label{ia:validation-and-evolution}

\begin{lemma}[Exact algebraic validity]\label{ia:exact-validator}
For each fixed $k$ there is a fixed finite mixture of bounded-probe
tests with perfect completeness on canonical extended labels and a
uniform constant gap on genuine degree-at-most-$D$ polynomial
families.  Passing all identities is equivalent to being the unique
canonical tuple of a valid extended label.  The gap is independent
of $L$ past the fixed cutoff and uniform over the allowed circuit
coefficients.  Linear probes have bounded point and mask density.
Every unknown product has compatible tags and can be queried with
independent full point marginals for its two operands.
\end{lemma}

\begin{proof}
The tests just constructed constitute a fixed number of family
types for fixed $k$; growing slot lists are checked by whole
contractions.  We prove exact soundness in an order that licenses
all substitutions.  First consider ignored-coordinate requirements,
mode and marker constancy, and the independence and coefficient
linearity devices of every linear tool.  If one fails, its direct
test has constant gap.  Otherwise these are exact polynomial
identities.  Monomial factorization and the nonrotating base calls
then force a Boolean tape, all its monomial extensions, its
union-of-stages validity, its output chart, and the canonical
per-primary certificates.  Mode and marker membership give unique
$\eta$ and $t$.  The legacy identity \eqref{ia:clock-legacy}
forces the unique frame $c$ and fixes $J_c$; clock-monomial
factorizations fix all $T_W$.

Now apply the clock-formula calls in construction order.  Their
uniqueness fixes each critical output, each derivative, and every
auxiliary.  A derivative's rows were computed from a prescribed
parent recipe; the proof never treats an unverified parent table
as an additional oracle for those rows.  The strictly decreasing
tier construction is therefore acyclic.  Finally the triangular
stack identities give \eqref{ia:stack-formula}.  Hence exactness
of all identities gives precisely a valid label's canonical tuple.
Conversely every valid label satisfies every identity for every
random choice in the support.

If exactness fails, the first failed identity in this ordering has
point degree $O(D+\ell h)$ after the stated preliminary devices
and bounded degree in its mask variables.  The polynomial zero and
contraction bounds give a uniform positive rejection probability.
Taking the minimum over the fixed mixture gives the asserted gap.
Honest degrees are $O(\ell h)<D$ eventually, also after any later
fixed-factor enlargement of the point domain.  The constants in
these degree bounds and in the finite family list depend on fixed
templates, not on the values of public circuit coefficients.

For the product claim, the only possibilities are monomial
factorizations, $J_cW$ factorizations, marker membership, and marker
times smaller stack.  In a factorization, choose the two native
operand arguments independently, assemble the output's relevant
groups from them, and independently fill all ignored coordinates.
Even repeated operand types then have independent full operand
points.  The relevant output groups are independent uniform too.
For marker products use independent points, justified by the
earlier constancy step.  Every such pair has compatible tags by
construction.  Bitness, diagonal grid entries, copy equalities,
and all other correlated grid uses are linear matrix rows, so they
require no concentrated unknown-product query.  The linear tool's
packing, permutation, and splitting probes have bounded individual
point density.  Dependent or prescribed masks, including those
inside products, are smoothed by fresh independent axis splittings.
These observations prove the stated probe interface.
\end{proof}

\begin{lemma}[All neighbor coordinates evolve by bounded affine pulls]
\label{ia:evolution}
For canonical valid data, edge existence and the transition branch
in either direction are determined using only mode, empty/full
marker comparisons, and boundedly many critical projected bits at
full progress.  Equality decisions can have any prescribed fixed
error probability by a fixed number of repetitions.  Once the
finite Boolean branch is fixed, every requested contraction in the
neighbor tuple is a bounded public linear sum of contractions in
the current tuple at invertible $K$-affine pullbacks of the common
point.  Its coefficients and slot pulls depend on the branch,
requested masks, and public data, but not on the point, row index,
or further hidden values.  Tag-$0$ updates use only tag $0$;
tag-$i$ updates use only compatible types.
\end{lemma}

\begin{proof}
We give formulas for each transition, including the boundaries.
Write $C^{S,t}_{F,\mathbf j}$ for the stack at progress $t$.

\paragraph{Progress increase and decrease.}
On increasing from $0$ to $1$, initialize every nonempty stack by
\begin{equation}
 C^{S,1}_{F,\mathbf j}
       =\mathbf1_{\mathbf j=\mathbf0}F,
 \qquad M^{\rm mark,new}_j=\mathbf1_{j=0}.         \label{ia:increase-zero}
\end{equation}
For $0<t<q-1$, and $\mathbf j$ in the padded radix range, let
$Z(\mathbf j)=\{\alpha\in S:j_\alpha=0\}$ and put
$k_\alpha=0$ for $\alpha\in Z(\mathbf j)$ and
$k_\alpha=j_\alpha-1$ otherwise.  Then
\begin{equation}
 C^{S,t+1}_{F,\mathbf j}
    =\left(\prod_{\alpha\in Z(\mathbf j)}D_\alpha\right)
           C^{S,t}_{F,\mathbf k},\qquad
 M^{\rm mark,new}_j=M^{\rm mark,old}_{j-1},       \label{ia:increase-positive}
\end{equation}
with zero padding.  For $j_\alpha>0$, the old exponent is
$t-1-(j_\alpha-1)=t-j_\alpha$, as required; for $j_\alpha=0$,
one extra dilation gives exponent $t$.  Old inactive coefficients
remain zero.  Partitioning by $Z(\mathbf j)$ and using the radix
pull formulas implements this with boundedly many mask operations.
For a decrease from $t>0$, the formula is simply
\begin{equation}
 C^{S,t-1}_{F,\mathbf j}
     =C^{S,t}_{F,(j_\alpha+1)_{\alpha\in S}},\qquad
 M^{\rm mark,new}_j=M^{\rm mark,old}_{j+1}.       \label{ia:decrease}
\end{equation}
Its exponent is $t-1-(j_\alpha+1)=t-2-j_\alpha$ and its support
is $j_\alpha<t-1$.  In particular $t=1$ produces empty stacks
and a zero marker.  Persistent data remain fixed during these
progress transitions.

\paragraph{Critical readings and full-progress edits.}
At $t=q-1$, for a critical chart $F$ with projection axis $\alpha$,
\begin{equation}
 \sum_{j=0}^{N-1}C^{\{\alpha\}}_{F,j}
     =\sum_{j=0}^{q-2}D_\alpha^{q-2-j}F
     =F\big|_{Z_\alpha=1}.                       \label{ia:critical-projection}
\end{equation}
All-ones radix masks perform this sum.  A critical output ignores
every other native group, so the result is a constant-in-point
logical or public rule bit.  Lemma~\ref{ia:validity} therefore
determines the edit mask from finitely many such bits.

Fix that mask and direction.  For every positive-tier parent $F$
and every $S\subseteq A_F$, including $S=\varnothing$, add the
increment
\begin{equation}
 \Delta C^S_{F,\mathbf j}
    =\sum_{r=0}^{N-1}
      C^{S\cup\{\alpha_{\rm new}\}}_{G_F^\sigma,(\mathbf j,r)}.
                                                               \label{ia:edit-increment}
\end{equation}
For a positive port parent use its corresponding legacy mask
$\chi$.  For active old indices, commuting all inherited-label
dilations with projection on the new label gives
\begin{align}
 \Delta C^S_{F,\mathbf j}
   &=\left(\prod_{\alpha\in S}D_\alpha^{q-2-j_\alpha}\right)
          G_F^\sigma(c,b;1,Y)                    \label{ia:edit-calculation}\\
   &=\left(\prod_{\alpha\in S}D_\alpha^{q-2-j_\alpha}\right)
          \bigl(F(c,b+h_\sigma(c);Y)-F(c,b;Y)\bigr).
                                                        \nonumber
\end{align}
If any inherited index is inactive, every summand is zero.  This
is precisely the new stack minus the old one, including the empty
stack, which is the persistent chart itself.  The recursion supplied
the necessary lower-tier derivative for \emph{every} positive-tier
parent, including proof and derivative charts, so all coordinates
update simultaneously from old data.  Tier-zero charts stay fixed
under the Boolean change.  The marker is unchanged and the mode
changes as prescribed.  Reverse edits obey the same calculation
with the inverse branch's coefficients.

\paragraph{Frame change and edge flags.}
At empty progress all stacks are zero.  For frame advance, a
rotating persistent chart is pulled back on its one designated
variable by multiplication by $\gamma$; for reverse advance use
$\gamma^{-1}$.  A nonrotating persistent chart stays fixed.  The
empty stacks remain zero, and the mode changes.  In particular,
an auxiliary with native variables $(R,I,X,U)$ rotates only $R$;
its packing, division, and coefficient formulas are not recomputed
by the update.

The successor rules are: pre increases unless full; at full pre
perform the enabled edit, or report the absent wrap edge; post
decreases unless empty; at empty post advance the frame.  The
predecessor rules are: full post performs the inverse enabled edit,
or reports the absent wrap edge; other post reverses a decrease;
positive-progress pre reverses an increase; empty pre reverses the
frame advance.  Empty must also be distinguished for initialization
in \eqref{ia:increase-zero}.  These are all finite cases.  Comparing
the marker with the public zero tensor or the public full tensor
uses random separable masks and binary forms, with a constant
inequality-detection gap.  A fixed number of repetitions makes
either comparison as reliable as prescribed.  Mode and critical
constant bits are read with forms taking $1$ to $1$.  No progress
index scan or proposed-target validity check is needed for edge
flags on the valid inner path.

Finally, all point maps in these formulas are identities, full-label
dilations \eqref{ia:dilation}, or full-label multiplications by
$\gamma^{\pm1}$, and their bounded compositions.  They are
invertible $K$-affine maps.  In \eqref{ia:edit-increment}, the child
retains the parent's placements and has an already available extra
whole label; the sum projects that label to one and makes it
irrelevant to the target coordinate.  Thus a requested update has
the form
\begin{equation}
 L\cdot F^{\rm new}(Y)
    =P_L(Y)+\sum_{r=1}^{R_*}a_r
                         (L_r\cdot F_r)(T_rY),    \label{ia:affine-update}
\end{equation}
where $R_*$ is bounded, $P_L$ is publicly computable, $a_r$ and
$L_r$ are fixed by the requested masks and branch, and $T_r$ is
an invertible $K$-affine map.  No $a_r$ or slot pull depends on
$Y$ or on a later row index.  The inherited tags in
Lemma~\ref{ia:derivative-closure}, and the public fixed branch,
give the tag assertions.
\end{proof}

\begin{remark}[Native tests and encoded updates have different maps]
\label{ia:testing-versus-evolution}
The arbitrary $\mathbb F_q$ coordinate permutations and surjective
digit packing in Lemma~\ref{ia:linear-tool} need not be $K$-affine.
They are used to construct certificates and to test systematic
target-point values with bounded density.  They are not maps for
propagating an encoded parity row or a line-coefficient array.
Equation~\eqref{ia:affine-update} is the separate evolution
interface for those encodings.  For example, the proof chart
$F(c,b;R,I,X,U)$ changes frame only through its whole $R$ label,
and changes Boolean tape through its stored child's full-label
projection.  Its canonical division is not rerun on the requested
encoded symbol.  This distinction permits the binary construction
to use ordinary systematic representations for all native tests
while transporting neighbor updates through its common affine
action and linear row encoding.
\end{remark}
\begin{lemma}[Source data, original indices, and Boolean degree]
\label{ia:source-data}
At each of the two labels in \eqref{ia:source-successor}, every
tag-$0$ chart and stack is publicly computable.  Every tag-$i$
chart and stack has exactly the value of its corresponding
positive-tag one-port-template chart at the same native point,
slot, and mask indices, evaluated on input $s_i$.  These statements
hold for the source and its first successor separately, with
their indicated progresses $q-1$ and $q-2$.

At either label, every prime-field coordinate function of every
chart value or coordinate coefficient has multilinear degree at
most three in the read inputs.  Any common public linear
representation of these arrays preserves this bound and their
original-index correspondence.  A binary bilinear form of two
such represented values has degree at most six.  Its coefficients
are computable in polynomial time by bounded-degree Boolean
interpolation.
\end{lemma}

\begin{proof}
At both labels, $c=\gamma^{-1}$, work is blank, $e_1=1$ and other
legacy bits are zero, the mode is post, and progress is the public
value specified in \eqref{ia:source-successor}.  Tag-$0$ persistent
recipes depend only on these values and public circuit data, so
they are computable.  Their derivatives might evaluate a translated
tape with nonblank work, but every flip vector is then public and
the translated work is public too.  Such a tape need not be valid:
the canonical recipe is evaluated by its finite-array operations,
and never by solving for a valid circuit computation.  In
particular a work-side copy certificate uses its own work primary
and public matrix, not a stored sum involving $s_i$.  Recursively,
all its derivative and proof arrays remain public.  Stack
dilations at either fixed progress preserve this conclusion.

Positive persistent recipes and layouts agree by
Lemma~\ref{ia:port-ancestry}, and their placements agree by
Lemma~\ref{ia:placements}.  The projection-axis choices, inherited
labels, $q,a,B$, and radix slot conventions also agree.  Formula
\eqref{ia:stack-formula} therefore gives identical stacks at each
of the two progresses, including every descendant proof stack.
This proves a literal map of original algebraic indices, not
merely equality up to unspecified isomorphisms or comparable sizes.
Applying the same public row extension, restriction to an ordered
line, extraction of its coefficients, slot contraction, or basis
map to corresponding arrays preserves both values and indices.
Later balancing by repeated copies may use different multiplicities;
it does not alter this statement about original indices.

At a fixed frame every persistent value and coefficient has
assigned degree at most three in the Boolean tape.  At the present
labels all tape bits except the read inputs are fixed public bits.
Fixed-progress stacks are fixed affine pullbacks or zero.  Their
degree is thus at most three in the read inputs as well.  A public
linear operation over $\mathbb F_q$, or over any fixed public
extension field followed by prime-field coordinates, is
$\mathbb F_2$-linear on these functions.  It cannot increase
multilinear degree.  For two degree-three coordinate vectors, a
binary bilinear form is a sum of products of two such coordinate
functions and hence has degree at most six.  Copies do not alter
degree.  Evaluate the specified canonical bit on supports of size
at most six and apply \eqref{ia:boolean-inversion}; the number of
evaluations and each canonical finite-array computation are
polynomial.  This proves the coefficient-computation assertion.
\end{proof}

\subsection{The complete inner algebra interface}\label{ia:interface}

\begin{lemma}[Inner algebra interface]\label{lem:inner-algebra-interface}
Fix $c_{\rm allow}$.  The field and grid parameters
\eqref{ia:parameters}--\eqref{ia:radix-size}, the common domain
$K^u$, and the one-port positive family templates can be chosen
before the port count and the function.  For every fixed $k$ and
fixed circuit size factor, and all sufficiently large $L$, a
uniform polynomial-time construction has the following properties.
\begin{enumerate}
\item \textbf{Paths and endpoints.}
The graph on all port tuples is the graph of
Definition~\ref{ia:extended-path}.  On each fixed tuple $s$ it is
one directed path with source and first successor
\eqref{ia:source-successor}, no isolated vertices or self edges,
and a sink at full pre stage $M-1$.  The point-constant output
chart \eqref{ia:output-chart} equals $f(s)$ at the sink.

\item \textbf{Canonical polynomial arrays and their size.}
Each valid extended label has a unique canonical tuple
$\mathcal A(c,b,\eta,t)$ of $\mathbb F_q$-coefficient polynomials
on $K^u$.  Persistent native dimensions are at most
$C_{\rm dim}\ell$, with $C_{\rm dim}$ independent of $k$, and
each persistent family has at most one $O(\ell)$ slot axis.
Stack families have a bounded number of slot axes, each of length
$O(\ell)$.  Their honest coordinate degrees are $O(\ell h)<D$.
For fixed $k$, the number of families is fixed and every full
evaluation table, fully masked table, public matrix, and indexing
map has polynomial size and is polynomial-time constructible.
The layouts for fixed $k,L$ are independent of the circuit's
values and wiring.  The construction does not enumerate all
Boolean tapes or graph vertices.

\item \textbf{Exact validity with smooth binary-form probes.}
There is a fixed finite mixture of bounded-probe algebraic tests,
perfectly complete on canonical valid tuples.  For every tuple
of genuine degree-at-most-$D$ polynomials, either it is a
canonical valid tuple or the rejection probability is at least
$\theta_{\rm alg}>0$.  For fixed $k$ this gap and the probe density
bounds are uniform in $L$ and in the allowed circuit coefficients.
Each linear probe has bounded point and mask marginals.  Every
unknown product has compatible tags and independent full point
marginals for its two operands, also for repeated types.
The tests require only a bounded number of additive binary forms
and binary bilinear forms of hidden field values, followed by a
bounded number of Boolean operations.  They do not require reading
the full binary expansion of a hidden field element.

\item \textbf{Invalid tapes and finite derivative closure.}
Canonical recipes exist also for arbitrary Boolean $b$ at every
specified $c,\eta,t$, whether or not that label is valid.  They
retain their coordinate bounds, logical degree at most three,
per-primary dependencies, and designated frame covariance.
Every positive-tier persistent family, including every
verification auxiliary, has all its required lower-tier
derivatives.  The recursion terminates in at most three tier
drops.  A nonvalid tuple so obtained has the same exact rejection
gap; no alternative certificate can hide a failed row condition.

\item \textbf{Neighbor evolution.}
On valid data, both edge flags and finite transition branches use
only mode, empty/full marker comparisons, and boundedly many
critical projected bits at full progress.  At any desired fixed
reliability, these distinctions use boundedly many binary probes.
After a branch is fixed, every neighbor contraction has the
form \eqref{ia:affine-update} with boundedly many invertible
$K$-affine pulls of the common point.  Coefficients and slot
pulls are independent of point and row position.  Tag-$0$
updates use tag $0$ only; positive updates use compatible tags.
The formulas update all proof data simultaneously.  Their
fixed-branch algebraic identities remain valid on canonical
invalid Boolean inputs, without asserting that the resulting
proposed label is valid.

\item \textbf{Standardized ports and source access.}
Every positive-tag family has the literal per-call ancestry and
original-index correspondence of Lemma~\ref{ia:port-ancestry},
including all descendants, full-label placements, projection
axes, and stacks.  At the source and first successor, the
tag-$0$ part is public, and tag $i$ agrees with the one-port
template on $s_i$.  The source Boolean-degree bounds are three
under public linear representations and six after one binary
bilinear pairing, as in Lemma~\ref{ia:source-data}.  These
templates and bounds are fixed before actual $k$ or any later
correction accuracy is chosen.
\end{enumerate}
\end{lemma}

\begin{proof}
The path and endpoint assertions are
Lemmas~\ref{ia:validity} and~\ref{ia:path-endpoints}.  The linear
tool, clock-formula construction, derivative closure, and common
placements give the persistent arrays.  Lemma~\ref{ia:stacks}
gives their stacks.  At fixed $k$, all depths and family counts
are fixed, while all dimensions and total mask lengths are
$O(\ell)$.  Consequently every domain and mask space has size
$q^{O(\ell)}=\poly(n')=\poly(L)$.  Row and coefficient arrays
have the same polynomial bounds.  Their construction consists
of the explicit finite-array operations and bounded-degree
interpolation above.  In particular, increasing the number of
separate calls does not change the already chosen common $u$.

Lemma~\ref{ia:exact-validator} proves the exact gap and all point
and tag restrictions.  Here is its conversion to the asserted
binary-form interface.  A nonzero element of $\mathbb F_q$ is
detected by a uniformly random additive form
$\mu:\mathbb F_q\to\mathbb F_2$ with probability $1/2$.
Applying $\mu$ to a field equality turns its public linear terms
into additive forms, since
\[
 \mu\left(\sum_r a_r x_r\right)
      =\bigoplus_r(\mu\circ[a_r])(x_r),
 \qquad [a_r](x)=a_rx.
\]
For a product, the requested bit is the binary bilinear form
$(x,y)\mapsto\mu(xy)$.  Random annihilating forms handle mode
membership as already specified.  A bounded field test therefore
uses a bounded number of such bits and Boolean operations,
losing at most a fixed factor in its gap.  Prescribed masks and
forms can be split by fresh additive randomness.  Their public
descriptions have polynomial size; even the space of all binary
bilinear forms on $\mathbb F_q$ has cardinality
\[
 2^{(\log_2 q)^2}=\poly(n'),
\]
since $\log q=O(\log\log n')$.  A bounded sequence of all the
displayed random choices consequently has polynomial support.

The invalid-tape assertion follows from the per-primary
canonical formulas and the finite-difference construction, not
from successful final rows.  If a label is invalid, a validity
or frame/legacy condition fails and exact soundness rejects;
uniqueness of the division certificates prevents a different
proof table from canceling the failed condition.
Lemma~\ref{ia:evolution} proves all fixed-branch update formulas,
including for canonical invalid tapes.  Frame changes occur
only when stacks are empty, so no commutation between centered
dilations and frame scaling is needed.  Finally
Lemmas~\ref{ia:port-ancestry}, \ref{ia:placements}, and
\ref{ia:source-data} give the exact standardized source
correspondence and its Boolean-degree bounds.  This proves
every part of the interface.
\end{proof}
